% !TeX root = elegantnote-en.tex
\section{Ketones and Aldehydes}

\subsection{Carbonyl Compounds 20260409}

\begin{enumerate}
    \item carbonyl group (\chemfig{C=O})
    \item Common Classes of Carbonyl Compounds % pic
    \item Ketone: Two alkyl groups bonded to a carbonyl group.
    
    Aldehyde: One alkyl group and one hydrogen bonded to a carbonyl group.
\end{enumerate}

\subsection{Structure of the Carbonyl Group 20260409}

\begin{enumerate}
    \item The carbonyl carbon atom is \textit{sp}$^2$ hybridized and bonded to three other atoms.
    \item The double bond of the carbonyl group has a large dipole moment because oxygen is more electronegative than carbon, and the bonding electrons are not shared equally.
    \item Resonance forms. The first resonance form is more important. %pic
\end{enumerate}

\subsection{Nomenclature of Ketones and AldehydesNomenclature of Ketones and Aldehydes 20260409}

\subsubsection{IUPAC Names}

\begin{enumerate}
    \item Ketone
    \begin{enumerate}
        \item Number the longest chain that includes the carbonyl carbon from the end closest to the carbonyl group.
        \item Replace the final \textit{-e} in the alkane name with \textit{-one}.
    \end{enumerate}
    %pic
    \item Aldehydes
    \begin{enumerate}
        \item An aldehyde carbon is number 1.
        \item Replace the final \textit{-e} of the alkane name with \textit{-al}.
        \item If the aldehyde group is a substituent of a large unit (usually a ring), the suffix \textit{carbaldehyde} is used.
    \end{enumerate}
    cyclohexanecarbaldehyde, 2-hydroxycyclopentane-1-carbaldehyde %pic
    \item A ketone or aldehyde group can also be named as a substituent on a molecule with a higher priority functional group as its root.
    \begin{enumerate}
        \item A ketone or aldehyde carbonyl is named by the prefix \textit{oxo} if it is included as part of the longest chain in the root name.
        \item When an aldehyde group is a substituent and not part of the longest chain, it is named by the prefix formyl.
    \end{enumerate}
    %Priority of Functional Groups in Naming Organic Compounds %pic
\end{enumerate}

\subsubsection{Common Names}

\begin{enumerate}
    \item Ketone common names are formed by naming the two alkyl groups bonded to the carbonyl group.
    \item Some ketones have historical common names.
    \item Common names of aldehydes are derived from the common names of the corresponding carboxylic acids. Greek letters are used with common names of aldehydes to give the locations of substituents. %pic
\end{enumerate}

\subsection{Physical Properties of Ketones and Aldehydes 20260409}

\begin{enumerate}
    \item Polarization of the carbonyl group creates dipole–dipole attractions between the molecules of ketones and aldehydes, resulting in higher boiling points than for hydrocarbons and ethers of similar molecular weights.
    \item Ketones and aldehydes have no \chemfig{O-H} or \chemfig{N-H} bonds, however, so their molecules cannot form hydrogen bonds with each other. Their boiling points are therefore lower than those of alcohols of similar molecular weight.
    \item H bond: Traditional (D-H···A), C-H···Cl, M-H, C-H···Benzene.
    \item The higher the \textit{s} orbital ratio, the stronger the base.
    \item Ketones and aldehydes have lone pairs of electrons and can act as hydrogen bond acceptors with other compounds having or \chemfig{O-H} or \chemfig{N-H} bonds. Ketones and aldehydes are good solvents for polar hydroxylic substances such as alcohols.
    \item Formaldehyde and acetaldehyde are the most common aldehydes.
    \item Formaldehyde is a gas at room temperature, so it is often stored and used as a 40\% aqueous solution called \textit{formalin}.
    \item When dry formaldehyde is needed, it can be generated by heating one of its solid derivatives. Trioxane is a cyclic trimer, containing three formaldehyde units. Paraformaldehyde is a linear polymer, containing many formaldehyde units.
    \item Acetaldehyde boils near room temperature.
    \item Acetaldehyde is also used as a trimer (paraldehyde) and a tetramer (metaldehyde).
\end{enumerate}

\subsection{Spectroscopy of Ketones and Aldehydes 20260409}

\subsubsection{Infrared Spectra of Ketones and Aldehydes}

\begin{enumerate}
    \item The carbonyl  stretching vibrations of simple ketones occur around \qty{1710}{\per \centi \meter}, and simple aldehydes around \qty{1725}{\per \centi \meter}.
    \item In addition to the carbonyl absorption, an aldehyde shows a set of two low-frequency stretching absorptions around 2710  and 2810 cm -1.
\end{enumerate}

\subsubsection{Proton NMR Spectra of Ketones and Aldehydes}

\begin{enumerate}
    \item Aldehyde protons normally appear at chemical shifts between $ \delta \ 9 $ and $ \delta \ 10 $.
    \item The aldehyde proton’s absorption may be split ($ J = 1 $ to $ 5 $ \si{Hz}) if there are protons on the atom.
\end{enumerate}
\subsubsection{Carbon NMR Spectra of Ketones and Aldehydes}

\begin{enumerate}
    \item The carbonyl carbon atoms of aldehydes and ketones have chemical shifts around 200 ppm in the carbon NMR spectrum.
\end{enumerate}

\subsubsection{Mass Spectra of Ketones and Aldehydes}

\begin{enumerate}
    \item $ \alpha $ cleavage %ppt
    \item McLafferty Rearrangement of Ketones and Aldehydes: a cyclic intramolecular transfer of a hydrogen atom from the $ \gamma $ carbon to the carbonyl oxygen.
    \item $ \gamma $ hydrogen is needed. %pic
    \item The McLafferty rearrangement is a characteristic fragmentation of ketones and aldehydes as long as they have hydrogens. It is equivalent to a cleavage between the $ \alpha $ and $ beta $ carbon atoms, plus one mass unit for the hydrogen that is transferred.
    \item The fragment from the McLafferty rearrangement has an even-numbered mass.
    \item enol
\end{enumerate}

\subsubsection{Ultraviolet Spectra of Ketones and Aldehydes}

\begin{enumerate}
    \item The $ \pi \rightarrow \pi^* $ Transition
\end{enumerate}

\subsection{Industrial Importance of Ketones and Aldehydes 20260409}

\subsection{Review of Syntheses of Ketones and Aldehydes 20260414}

\subsubsection{Ketones and Aldehydes from Oxidation of Alcohols}

\begin{enumerate}
    \item Often use a Grignard reagent to synthesize an alcohol with the correct structure and oxidize it to the final product.
    \begin{enumerate}
        \item Secondary alcohols $ \rightarrow $ ketones: Sodium dichromate in sulfuric acid (“chromic acid,” \ce{H2CrO4}) is the traditional laboratory reagent for oxidizing secondary alcohols to ketones.
        \item Primary alcohols $ \rightarrow $ aldehydes: Pyridinium chlorochromate (PCC), a complex of chromium trioxide with pyridine and \ce{HCl}, provides good yields of aldehydes without over-oxidation.
    \end{enumerate}
\end{enumerate}

\subsubsection{Ketones and Aldehydes from Ozonolysis of Alkenes}

\begin{enumerate}
    \item Ozonolysis, followed by a mild reduction, cleaves alkenes to give ketones and aldehydes.
    \item Ozonolysis is useful as a synthetic method.
\end{enumerate}

\subsubsection{Phenyl Ketones and Aldehydes: Friedel–Crafts Acylation}

\begin{enumerate}
    \item Friedel–Crafts acylation is an excellent method for making alkyl aryl ketones or diaryl ketones.
\end{enumerate}

\subsubsection{Ketones and Aldehydes from Hydration of Alkynes}

\begin{enumerate}
    \item Catalyzed by Acid and Mercuric Salts
    \begin{enumerate}
        \item Hydration of a terminal alkyne, catalyzed by a combination of sulfuric acid and mercuric ion.
        \item The initial product of Markovnikov hydration is an enol.
    \end{enumerate}
    \item Hydroboration–Oxidation of Alkynes
    \begin{enumerate}
        \item Hydroboration–oxidation of an alkyne gives antiMarkovnikov addition of water across the triple bond.
        \item On oxidation of the borane, the unstable enol quickly tautomerizes to an aldehyde.
    \end{enumerate}
\end{enumerate}

\subsection{Synthesis of Ketones and Aldehydes Using 1,3-Dithianes 20260414}

\begin{enumerate}
    \item 1,3-Dithiane can be deprotonated by strong bases such as \textit{n}-butyllithium.
    \item Alkylation of the dithiane anion by a primary alkyl halide or a tosylate gives a \textbf{thioacetal} that can be hydrolyzed into the aldehyde by using \textbf{an acidic solution of mercuric chloride}. %pic
    \item The thioacetal can be isolated and deprotonated. Alkylation gives a \textbf{thioketal}, then hydrolysis will produce a ketone.
\end{enumerate}

\subsection{Synthesis of Ketones from Carboxylic Acids 20260414}

\begin{enumerate}
    \item Organolithium reagents can be used to synthesize ketones from carboxylic acids.
    \item Organolithiums will attack the lithium salts of carboxylate anions to give dianions.
    \item Protonation of the dianion forms the hydrate of a ketone, which quickly loses water to give the ketone.
\end{enumerate}

\subsection{Synthesis of Ketones and Aldehydes from Nitriles 20240414}

\begin{enumerate}
    \item Nitriles are compounds containing the cyano (\chemfig{-C~N}) functional group.
    \item A Grignard or organolithium reagent can attack the carbon of the nitrile.
    \item The imine is then hydrolyzed to form a ketone.
    \item Aluminum hydrides can reduce nitriles to the corresponding aldehydes. Diisobutylaluminum hydride, abbreviated (i-Bu)2AlH or \textbf{DIBAL-H}, is commonly used for the reduction of nitriles (mild reducing agent).
\end{enumerate}

\subsection{Synthesis of Aldehydes and Ketones from Acid Chlorides and Esters 20260414}

\begin{enumerate}
    \item Acids can be reduced to aldehydes by first converting them to acid chloride.
    \item Lithium tri-\textit{tert}-butoxyaluminum hydride is a milder reducing agent that reacts faster with acid chlorides than with aldehydes.
    \item Synthesis of Ketones
    \begin{enumerate}
        \item Grignard and organolithium reagents add to acid chlorides to give ketones, but they add again to the ketones to give tertiary alcohols.
        \item \textbf{Lithium dialkylcuprates} (Gilman reagent) and other organocuprates react with acid chlorides to give good yields of a wide variety of ketones.
        \item The lithium dialkylcuprate is formed by the reaction of two equivalents of the corresponding organolithium reagent with cuprous iodide.
    \end{enumerate}
    \item Syntheses of Ketones and Aldehydes %pic
\end{enumerate}

\subsection{Reactions of Ketones and Aldehydes: Introduction to Nucleophilic Addition 20260414}

\begin{enumerate}
    \item Strong Nu: \ce{R-}, \ce{H-}, \ce{NaNH2}, \ce{RO-}, \ce{NC-}, \ce{\chemfig{C~C}-}
    
    Weak Nu: \ce{ROH}, \ce{RNH2}, \ce{H2O}.
    \item Weak nucleophiles can add to activated carbonyl group under acidic conditions. A carbonyl group that is protonated (or bonded to some other electrophile) is strongly electrophilic, inviting attack by a weak nucleophile.
    \item Aldehydes are more reactive than ketones toward nucleophilic additions. The enhanced reactivity of aldehydes is due to an electronic effect and a steric effect.
\end{enumerate}

\begin{example} % 18-12
    Sodium triacetoxyborohydride, \ce{NaBH(OAc)3}, is a mild reducing agent that reduces aldehydes much more quickly than ketones. It can be used to reduce aldehydes in the presence of ketones, such as in the following reaction: %三乙酰氧基硼氢化钠
\end{example}

\begin{mechanism}
    Nucleophilic Additions to Carbonyl Groups
    \begin{enumerate}
        \item Basic Conditions (strong nucleophile)
        \begin{itemize}
            \item Step 1: A strong nucleophile adds to the carbonyl group to form an alkoxide.
            \item Step 2: A weak acid protonates the alkoxide to give the addition product.
        \end{itemize}
        \item Acidic Conditions (weak nucleophile, activated carbonyl)
        
        hemiacetal %半缩醛
    \end{enumerate}
\end{mechanism}

\subsection{The Wittig Reaction 20260414 20260416}

\begin{enumerate}
    \item Adding a phosphorus-stabilized carbanion to a ketone or aldehyde, the Wittig reaction converts the carbonyl group into a new \chemfig{C=C} double bond where no bond existed before.
    \item A phosphorus ylide is used as the nucleophile in the reaction. %pic
    \item Triphenylphosphine oxide is exceptionally stable, and the conversion of triphenylphosphine to triphenylphosphine oxide provides the driving force for the Wittig reaction.
    \item Mixtures of cis and trans isomers often result.
\end{enumerate}

\begin{mechanism}
    The Wittig Reaction
    \begin{itemize}
        \item Step 1: The ylide attacks the carbonyl to form a betaine.
        \item Step 2: The betaine closes to a four-membered ring oxaphosphetane (first \chemfig{P-O} bond formed). %oxa thia aza 氧杂 硫杂 氮杂
        \item Step 3: The ring collapses to the products (second \chemfig{P-O} bond formed).
    \end{itemize}
\end{mechanism}

\begin{example} %18-14
    Like other strong nucleophiles, triphenylphosphine attacks and opens epoxides. The initial product (a betaine) quickly cyclizes to an oxaphosphetane that collapses to an alkene and triphenylphosphine oxide.
    \begin{enumerate}[label=(\alph*)]
        \item Show each step in the reaction of \textit{trans}-2,3-epoxybutane with triphenylphosphine to give but-2-ene. What is the stereochemistry of the double bond in the product?
        \item Show how this sequence might be used to convert \textit{cis}-cyclooctene to \textit{trans}-cyclooctene.
    \end{enumerate}
\end{example}

\begin{enumerate}[resume]
    \item Planning a Wittig Synthesis
    \begin{enumerate}
        \item In general, the ylide should come from an unhindered alkyl halide.
    \end{enumerate}
\end{enumerate}

\begin{example} %18-16
    Show how Wittig reactions might be used to synthesize the following compounds. In each case, start with an alkyl halide and a ketone or an aldehyde.
    \begin{enumerate}[label=(\alph*)]
        \item \chemfig{Ph-CH=C(CH_3)_2}
        \item \chemfig{Ph-C(CH_3)=CH_2}
    \end{enumerate}
\end{example}

\subsection{Hydration of Ketones and Aldehydes 20260416}

\begin{enumerate}
    \item In an aqueous solution, a ketone or an aldehyde is in equilibrium with its hydrate, a geminal diol. %偕二醇 同碳二醇
    \item With most ketones, the equilibrium favors the unhydrated keto form of the carbonyl.
    \item Aldehydes are more likely than ketones to form stable hydrates, because an aldehyde carbonyl has only one stabilizing alkyl group. The partial positive charge of the aldehyde is not as well stabilized.
\end{enumerate}

\begin{mechanism}
    Hydration of Ketones and Aldehydes
    \begin{enumerate}
        \item \textbf{In acid}
        
        The acid-catalyzed hydration is a typical acid-catalyzed addition to the carbonyl group. Protonation, followed by addition of water, gives a protonated product. Deprotonation gives the hydrate.
        \begin{itemize}
            \item Step 1: Protonation.
            \item Step 2: Water adds.
            \item Step 3: Deprotonation.
        \end{itemize}
    \end{enumerate}
\end{mechanism}

\subsection{Formation of Cyanohydrins 20260416}

\begin{enumerate}
    \item Cyanide ion is a strong base and a strong nucleophile. It attacks ketones and aldehydes to give addition products called cyanohydrins.
    \item These equilibrium constants follow the general reactivity trend of ketones and aldehydes: 
    \begin{equation*}
        \text{formaldehyde} > \text{other aldehydes} > \text{ketones}
    \end{equation*}
\end{enumerate}

\begin{mechanism}
    Formation of Cyanohydrins

    Cyanohydrin formation is a perfect example of base-catalyzed addition to a carbonyl group. The strong nucleophile adds in the first step to give an alkoxide. Protonation gives the cyanohydrin.
    \begin{itemize}
        \item Step 1: Cyanide adds to the carbonyl.
        \item Step 2: Protonation gives the cyanohydrin.
    \end{itemize}
\end{mechanism}

\begin{enumerate}[resume]
    \item Organic compounds containing the cyano group are called nitriles. A cyanohydrin is therefore an $\alpha$-hydroxynitrile.
    \item Nitriles hydrolyze to carboxylic acids under acidic conditions, so cyanohydrins hydrolyze to $\alpha$-hydroxy acids.
\end{enumerate}

\subsection{Formation of Imines 20260416}

\begin{enumerate}
    \item Ammonia or a primary amine reacts with a ketone or an aldehyde to form an imine.
    \item Imines are nitrogen analogues of ketones and aldehydes, with a \chemfig{C=N} in place of the carbonyl group.
    \item Imines are basic, a substituted imine is also called a \textbf{Schiff base}. %手性化合物的前驱体
    \item The proper pH is crucial to imine formation, fastest around pH 4.5.
\end{enumerate}

\begin{mechanism}
    Formation of Imines
    \begin{enumerate}
        \item First part: Acid-catalyzed addition of the amine to the carbonyl group.
        \begin{itemize}
            \item Step 1: Protonation of the carbonyl.
            \item Step 2: Addition of the amine.
            \item Step 3: Deprotonation.
        \end{itemize}
        \item Second part: Acid-catalyzed dehydration.
        \begin{itemize}
            \item Step 4: Protonation of the \chemfig{-OH} group.
            \item Step 5: Loss of \ce{H2O}.
            \item Step 6: Deprotonation.
        \end{itemize}
    \end{enumerate}
\end{mechanism}

\begin{enumerate}[resume]
    \item Imine formation is reversible, and most imines can be hydrolyzed back to the amine and the ketone or aldehyde. %pic
    \item Enamine Formation %烯胺
    
    An enamine results from the reaction of a ketone or aldehyde with a secondary amine. Activate the $\alpha$ place of carbonyl group.
\end{enumerate}

\subsection{Condensations with Hydroxylamine and Hydrazines 20260416}

\begin{enumerate}
    \item Condensations of Amines with Ketones and Aldehydes %pic
    \item oxime %肟 
    hydrazine %肼
    hydrazone %腙
    phenylhydrazine %苯肼
    phenylhydrazone %苯腙
    semicarbazide %半胱氨基脲
    semicarbazone %半胱氨基脲腙
\end{enumerate}

\begin{example} % 18-25
    2,4-Dinitrophenylhydrazine is frequently used for making derivatives of ketones and aldehydes because the products (2,4-dinitrophenylhydrazones, called 2,4-DNP derivatives) are even more likely than the phenylhydrazones to be solids with sharp melting points. Propose a mechanism for the reaction of acetone with 2,4-dinitrophenylhydrazine in a mildly acidic solution. 
    
    It is used in test for glucose. %2,4-二硝基苯肼
\end{example}

\begin{example}
    Predict the products of the following reactions.
    \begin{enumerate}[label=(\alph*)]
        \item
    \end{enumerate}
\end{example}

\subsection{Formation of Acetals 20260416}

\begin{enumerate}
    \item Ketones and aldehydes react with alcohols to form acetals.
    \item In the formation of an acetal, two molecules of alcohol add to the carbonyl group, and one molecule of water is eliminated.
    \item Acetal formation must be acidcatalyzed.
\end{enumerate}

\begin{mechanism}
    Formation of Acetals
    \begin{enumerate}
        \item First half: Acid-catalyzed addition of the alcohol to the carbonyl group.
        \begin{itemize}
            \item Step 1: Protonation.
            \item Step 2: Alcohol adds.
            \item Step 3: Deprotonation.
        \end{itemize}
        \item Second half: $\mathrm{S_N}1$ substitution of the protonated hemiacetal.
        \begin{itemize}
            \item Step 4: Protonation.
            \item Step 5: Loss of water.
        \end{itemize}
    \end{enumerate}
\end{mechanism}

\begin{enumerate}[resume]
    \item Equilibrium of Acetal Formation
    \item Cyclic Acetala
    \begin{enumerate}
        \item Formation of an acetal using a diol as the alcohol gives a cyclic acetal.
        \item Ethylene glycol is often used to make cyclic acetals; its acetals are called ethylene acetals (or ethylene ketals).
        \item The reaction is reversible and is used in synthesis to protect carbonyls from reaction
    \end{enumerate}
\end{enumerate}

\subsection{Use of Acetals as Protecting Groups 20260416 20260421}

\begin{example} % 18-33
    Show how you would accomplish the following syntheses. You may use whatever additional reagents you need.
\end{example}

\subsection{Oxidation of Aldehydes 202640416}

\begin{enumerate}
    \item Aldehydes are easily oxidized to carboxylic acids by common oxidants such as bleach (sodium hypochlorite), chromic acid, permanganate, and peroxy acids.
    \item Silver ion oxidizes aldehydes selectively in a convenient functional-group test for aldehydes.
    \item The \textbf{Tollens test} involves adding a solution of silver–ammonia complex (the \textbf{Tollens reagent}) to the unknown compound.
\end{enumerate}

\begin{example} % 18-65
    f 
\end{example}

\subsection{Reductions of Ketones and Aldehydes 20250421}

\subsubsection{Hydride Reductions (Review)}

\subsubsection{Catalytic Hydrogenation}

\subsubsection{Deoxygenation of Ketones and Aldehydes}

\begin{enumerate}
    \item Clemmensen Reduction (Review)
    \item Wolff–Kishner Reduction: The ketone or aldehyde is converted to its hydrazone, which is heated with a strong base such as KOH or potassium tert-butoxide. Compounds that cannot survive treatment with hot acid can be deoxygenated using the Wolff–Kishner reduction.
\end{enumerate}

Assignment 18-39, 48, 50/abcd, 56, 60, 61

\begin{example} % 18-40
    Rank the following carbonyl compounds in order of increasing equilibrium constant for hydration:

    \ce{CH3COCH3Cl} \quad \ce{ClCH2CHO} \quad \ce{CH2O} \quad \ce{CH3COCH3} \quad \ce{CH3CHO}
\end{example}

\begin{example} % 18-44
    The proton NMR spectrum of a compound of formula \ce{C10H12O} follows. This compound reacts with an acidic solution of 2,4-dinitrophenylhydrazine to give a crystalline derivative, but it gives a negative Tollens test. Propose a structure for this compound and give peak assignments to account for the signals in the spectrum.
\end{example}

\begin{example} % 18-66 *
    Solving the following road-map problem depends on determining the structure of \textbf{A}, the key intermediate. Give structures for compounds \textbf{A} through \textbf{K}.
\end{example}

半缩醛在托伦试剂是否稳定

氧化银氧化醛、醇有选择性

康尼扎罗反应 

安息香缩合

\begin{enumerate}
    \item Benzoin Condensation
    \begin{figure}[h]
        \centering
        \includegraphics[width=0.7\textwidth]{18.21_Benzoin_Condensation.png}
        \caption{Benzoin Condensation}
        \label{18.21_Benzoin_Condensation}
    \end{figure}

\end{enumerate}